\section{The reflected form on an adapted torus}\label{sec:finite-reflection}

The chamber bonds crossing a mirror join reflected endpoints. Their
Clifford projections give positive fermion coefficients. We combine
this expansion with a Fourier representation of the integer closedness
constraints to obtain a positive form, including all flux sectors.
This is the adapted-torus reflected form of
\cite[Sections~3.6--3.7]{ChamberQED}, whose two seam contributions
we retain explicitly.

\subsection{Symmetries and reflected Grassmann products}
\label{subsec:reg-symmetries}

The field strength transforms as a two-cochain, with the orientation
of each plaquette included. For an antiperiodic cubic torus the law is invariant
under $W$ and under $F\mapsto-F$. We record the corresponding
fermionic identities, since reflection positivity uses their signs.

Gauge covariance follows directly from
\[
 U_{A-d\lambda}(x,x')
 =e^{iq\lambda(\pi(x))}U_A(x,x')e^{-iq\lambda(\pi(x'))}.
\]
For lattice covariance set
$\gamma_5=\left(\begin{smallmatrix}0&1_2\\1_2&0\end{smallmatrix}\right)$.
It anticommutes with the Clifford triple, and
\[
 (\gamma_n\gamma_5)\gamma_v(\gamma_n\gamma_5)^{-1}
 =\gamma_{r_nv}.
\]
If the linear part of $g\in W$ is $r_{n_1}\cdots r_{n_k}$, the
corresponding product of these spin matrices, followed by the
permutation $x\mapsto g^{-1}x$, intertwines $\mathcal D_A^{\mathbb T}$ and
$\mathcal D_{(g^{-1})^*A}^{g\mathbb T}$, where
$g\mathbb T=(r\Gamma,\chi\circ r^{-1})$ and $r$ is the linear part
of $g$. Taking determinants and changing
variables in \eqref{eq:reg-functional} proves
\[
 \mu_{g\mathbb T}=((g^{-1})^*)_*\mu_{\mathbb T}.
\]
The configuration measure is preserved because pullback permutes
cochain coordinates up to signs and maps one fundamental domain
to another. All translations preserve $\mathbb T$, and all elements
of $W$ preserve an antiperiodic cubic torus, proving the stated
invariance.

For charge conjugation let $C=\sigma_2\oplus\sigma_2$. The Pauli
identity $C\gamma_v^\top C^{-1}=-\gamma_v$, together with reversal
of the oriented bond, gives
\[
 C(\mathcal D_A^{\mathbb T})^\top C^{-1}=\mathcal D_{-A}^{\mathbb T}.
\]
Thus the determinant is unchanged by $A\mapsto-A$. The substitution
$(A,z)\mapsto(-A,-z)$ preserves the configuration measure and sends
$F$ to $-F$, proving charge symmetry of the curvature law.

At a mirror $\Pi$ with upward normal $n$, the Grassmann reflection
$\Theta_\Pi$ is antilinear and reverses products. In components it is
\begin{equation}\label{eq:reg-grassmann-reflection}
 \Theta_\Pi\psi_\alpha(x)
  =-\sum_\beta\bar\psi_\beta(\theta_\Pi x)(\gamma_n)_{\beta\alpha},
 \qquad
 \Theta_\Pi\bar\psi_\alpha(x)
  =-\sum_\beta(\gamma_n)_{\alpha\beta}\psi_\beta(\theta_\Pi x).
\end{equation}
On scalar coefficient functions it acts as
$\Phi(A,z)\mapsto\overline{\Phi(\theta_\Pi^*A,\theta_\Pi^*z)}$.
The Clifford identities give $\Theta_\Pi^2=1$ and
\begin{equation}\label{eq:reg-reflected-bilinear}
 \Theta_\Pi(\bar\psi(x)Y\psi(x'))
 =\bar\psi(\theta_\Pi x')\gamma_nY^*\gamma_n\psi(\theta_\Pi x).
\end{equation}
For a hopping term,
$\gamma_nw(x,x')\gamma_n=w(\theta_\Pi x',\theta_\Pi x)$, by
\eqref{eq:reg-clifford-reflection}. Reflection therefore permutes the
oriented terms of the action, including the real mass terms.

If $Y$ is a finite set of sites, then
\begin{equation}\label{eq:reg-reflected-berezin}
 \int_{\theta_\Pi Y}\Theta_\Pi\xi=\overline{\int_Y\xi}.
\end{equation}
Only the top-degree coefficient needs to be checked. At one site
reflection maps $\prod_\alpha\psi_\alpha\bar\psi_\alpha$ to the same
product at the reflected site, multiplied by $(\det\gamma_n)^2=1$.
Products from different sites are even. This proves
\eqref{eq:reg-reflected-berezin}, and a measure-preserving change of
photon variables gives
\begin{equation}\label{eq:reg-reflected-functional}
 Z_{\theta_\Pi\mathbb T}[\Theta_\Pi G]
 =\overline{Z_{\mathbb T}[G]}
\end{equation}
for bounded coefficient integrands, and for the integrable
Gaussian-growth classes considered below.

On even integrands, compositions of these reflections define an
action of $W$. The ambiguity of a reflection word acts only by the
common sign on all fermion generators. To verify this, compose two
words representing the same lattice isometry. The resulting even
word has trivial geometric action, and its spin matrix is
$\diag(u,u)$ with $u$ a product of an even number of Pauli matrices
$\sigma_n$. It commutes with all Pauli matrices, hence is scalar;
its determinant is one, so $u=\pm1_2$. Its action on an even monomial
is consequently the identity. We write $G^g$ for this action when
folding even integrands. For a real scalar observable it is
$G^g(F)=G(g^*F)$.

\subsection{The fermion crossing at a mirror}\label{subsec:reg-crossing}

Let a bond of type $j$ cross $\Pi$, with lower endpoint $x$ and upper
endpoint $x'=\theta_\Pi x$. Its direction is the upward normal $n$.
The two hopping matrices are
\begin{equation}\label{eq:reg-crossing-projectors}
 \begin{split}
 w(x,x')&=\tfrac{w_j}{2}\bigl((1+\rho)P_n^--(1-\rho)P_n^+\bigr),\\
 w(x',x)&=\tfrac{w_j}{2}\bigl((1+\rho)P_n^+-(1-\rho)P_n^-\bigr).
 \end{split}
\end{equation}
Choose orthonormal bases $e_{n,1}^\pm,e_{n,2}^\pm$ of the two
projected spaces. For each $k=1,2$ introduce four crossing modes.
Their coefficients and phases are displayed together:
\begin{equation}\label{eq:reg-crossing-modes}
\begin{array}{c|c|c|c}
 &\zeta&\omega&u\\ \hline
 1&(e_{n,k}^-)^*\psi(x')&\tfrac{w_j}{2}(1+\rho)&U_A(x,x')\\
 2&(e_{n,k}^+)^*\psi(x')&\tfrac{w_j}{2}(1-\rho)&U_A(x,x')\\
 3&\bar\psi(x')e_{n,k}^+&\tfrac{w_j}{2}(1+\rho)&\overline{U_A(x,x')}\\
 4&\bar\psi(x')e_{n,k}^-&\tfrac{w_j}{2}(1-\rho)&\overline{U_A(x,x')}
\end{array}
\end{equation}
All coefficients are positive in the range $0<\rho<1$.

\begin{proposition}\label{prop:reg-crossing-expansion}
Let $J$ be a finite collection of bonds crossing $\Pi$, and let
$S_\Pi^\times(J)$ be their part of the action. Index all their crossing
modes by a finite set $\mathcal I$. Then
\begin{equation}\label{eq:reg-crossing-positive}
 -S_\Pi^\times(J)=\sum_{i\in\mathcal I}\omega_i u_i\Theta_\Pi(\zeta_i)\zeta_i.
\end{equation}
Choose an order on $\mathcal I$ and, for $I\subset\mathcal I$, put
$\zeta_I=\prod_{i\in I}\zeta_i$, $\omega_I=\prod_{i\in I}\omega_i$
and $u_I=\prod_{i\in I}u_i$. Then
\begin{equation}\label{eq:reg-crossing-exponential}
 e^{-S_\Pi^\times(J)}=\sum_{I\subset\mathcal I}\omega_Iu_I\Theta_\Pi(\zeta_I)\zeta_I.
\end{equation}
At site and diagonal mirrors all $u_i$ are one. At a link plane,
groups one and two have charge $+1$, groups three and four charge
$-1$, and their phases are $e^{iq\mathfrak q(i)A(b)}$ at the cut
photon link $b$.
\end{proposition}

\begin{proof}
The reflections of the four modes in \eqref{eq:reg-crossing-modes}
are, respectively,
\[
 \bar\psi(x)e_{n,k}^-,\quad-\bar\psi(x)e_{n,k}^+,
 \quad-(e_{n,k}^+)^*\psi(x),\quad(e_{n,k}^-)^*\psi(x).
\]
Multiplying each reflected mode by the original mode and using
anticommutation gives the four projected terms in
\eqref{eq:reg-crossing-projectors}, with the indicated phases.
Summation over $k$ uses
$P_n^\pm=\sum_ke_{n,k}^\pm(e_{n,k}^\pm)^*$ and gives the negative of
the two bond terms. Summation over bonds proves
\eqref{eq:reg-crossing-positive}.

Each $\Theta_\Pi(\zeta_i)\zeta_i$ is even and has square zero.
These elements commute. Reversal of the order under $\Theta_\Pi$
gives
\[
 \prod_{i\in I}\Theta_\Pi(\zeta_i)\zeta_i
 =\Theta_\Pi(\zeta_I)\zeta_I,
\]
with no remaining sign. Expanding the product of the finite
exponentials proves \eqref{eq:reg-crossing-exponential}.
Only type-$3$ bonds carry a photon phase, and they cross link
planes. Their projected endpoints are the oriented endpoints of
the cut link, giving the last assertion.
\end{proof}

\subsection{Adapted tori and the second mirror}\label{subsec:reg-adapted}

An adapted torus for $\ell\in\mathcal L$ has the form
\begin{equation}\label{eq:reg-adapted-torus}
 \Gamma=\Gamma_\parallel\oplus\Z v,
 \quad \Gamma_\parallel=\Gamma\cap\{\ell=0\},
 \quad v=\frac{p}{|n_\ell|^2}n_\ell\in\Z^3,
 \quad \chi(v)=-1,
\end{equation}
where $p=\ell(v)>0$ is the height period. For an admissible $c$ set
\[
 c'=c+p/2,\quad \Pi=(\ell,c),\quad\Pi'=(\ell,c'),\quad
 U_+=\{c<\ell<c'\},\quad U_-=\{c-p/2<\ell<c\}.
\]
The two mirrors are admissible: for coordinate families $p\ge4$,
while for diagonal families $p$ is even and $p\ge8$. Their
reflections satisfy $\theta_{\Pi'}=\tau_v\theta_\Pi$ and induce the
same geometric reflection on the torus, exchanging its two halves.
An antiperiodic coordinate torus is adapted to each coordinate
family. The torus with periods $ae_1-ae_2$, $be_1+be_2$ and $ce_3$,
antiperiodic on these periods, is adapted to the corresponding two
diagonal families and the third coordinate family. Cubic images
give adapted tori for all nine families.

Choose representatives $\mathcal X_+$ in $U_+$ and
$\mathcal X_-=\theta_\Pi\mathcal X_+$ in $U_-$. They partition the
fermion sites of the torus. Let $S_+$ contain the mass terms and
oriented bonds entirely in $U_+$, and set
$\theta=\theta_\Pi$, $\Theta=\Theta_\Pi$; the second boundary mirror
induces the same geometric reflection on the torus.
The reflected part is $\Theta S_+$. Its coefficient functions
involve only links in the closed upper slab.

The crossing modes at $\Pi$ lie in
$\operatorname{Gr}(\mathcal X_+)$. At $\Pi'$ their upper endpoints
are outside $U_+$; write
$\check\zeta_{i'}=\Theta_{\Pi'}\zeta_{i'}$, so the checked modes
lie in $\operatorname{Gr}(\mathcal X_+)$. The boundary condition gives
\begin{equation}\label{eq:reg-second-reflection}
 \Theta_{\Pi'}=\chi(v)^{\deg}\Theta_\Pi
                 \quad\hbox{on }\operatorname{Gr}_{\mathbb T}.
\end{equation}
In the crossing term at $\Pi'$ we therefore have
\[
 \Theta_{\Pi'}(\zeta_{i'})\zeta_{i'}
 =-\chi(v)\Theta(\check\zeta_{i'})\check\zeta_{i'}.
\]
The minus sign is from exchanging the two odd factors. Since
$\chi(v)=-1$, it is canceled by antiperiodicity. Both boundary
crossings thus have the same positive coefficients, and
\begin{equation}\label{eq:reg-two-crossings}
 e^{-S_{\mathbb T}}=e^{-S_+}\Theta(e^{-S_+})
 \sum_{I,I'}\omega_I\omega_{I'}u_Iu_{I'}
       \Theta(\zeta_I)\zeta_I
       \Theta(\check\zeta_{I'})\check\zeta_{I'}.
\end{equation}
Here $I$ and $I'$ range over the finite crossing-mode sets at the
two mirrors.

The reflected Berezin integral separates the two halves:
\begin{equation}\label{eq:reg-half-berezin}
 \int_{\mathbb T}\Theta(\xi')\xi
 =\overline{\int_{\mathcal X_+}\xi'}
                  \int_{\mathcal X_+}\xi,
 \qquad \xi,\xi'\in\operatorname{Gr}(\mathcal X_+).
\end{equation}
Indeed, the top monomial is the product of the two half-top
monomials, and reflection maps one to the other by
\eqref{eq:reg-reflected-berezin}. We must still separate the photon
flux constraints; these couple the halves at the cut cubes.

\subsection{The photon crossing and Gauss' law}\label{subsec:reg-gauss}

Let $\mathcal M$ be the two boundary mirrors, understood periodically.
Classify a cell by its relative interior: it is upper or lower if
that interior is in the respective half, lies on the mirrors if the
whole cell does, and is cut if its relative interior meets a mirror
without being contained in it. Write $B_+,B_-,B_0,B_\times$ for the
four sets of links, $P_+,P_-,P_0,P_\times$ for the plaquettes, and
$Q_+,Q_-,Q_\times$ for the cubes.

Reflection exchanges the upper and lower cells. It fixes a cell
lying on the mirrors with orientation sign $+1$, and a cut cell
with orientation sign $-1$. Each cell meets at most one boundary
mirror: its height variation is at most $|n_\ell|^2$, smaller than
$p/2$. To justify the orientation assertion, a point in the
intersection of the relative interior with the mirror is fixed
and lies in the relative interior of the reflected cell. Relative
interiors of different cells are disjoint, so the cell is fixed.
Reflection is the identity on its affine hull in the first case
and a reflection of that affine hull in the cut case.

An upper cell has only upper faces or faces lying on the mirrors;
a cut cell has only upper, lower or cut faces. Moreover a plaquette
containing a cut link is cut, and a cube containing a cut plaquette
is cut. An upper face can touch the boundary at a relative interior
point only when its affine height is constant. A face of a cut cell
cannot lie in the mirror: reflection in the supporting plane of a
face would not preserve that cell. Finally, the height takes both
signs on a cut face, hence on the relative interior of its coface.
In particular, the incidences at a cut plaquette and a cut cube give
\begin{equation}\label{eq:reg-cut-differentials}
 (dA)_\sigma=(d_+A)_\sigma-(d_+\theta^*A)_\sigma
                                      +(d_\times A)_\sigma,
 \qquad
 (dz)(Q)=(d_+z)(Q)-(d_+\theta^*z)(Q)+(d_\times z)(Q),
\end{equation}
where $d_+$ and $d_\times$ retain the incidences from upper and cut
faces, respectively. For example,
$(d_+A)_\sigma=\sum_{b\in B_+}[\sigma:b]A(b)$.
The minus sign in the lower term follows from orientation reversal
of the cut cell and the identity $d\theta^*=\theta^*d$.

Use representatives
\[
 A=(A_+,A_-,A_0,a),\qquad z=(z_+,z_-,z_0,z_\times).
\]
Choose the interval $[0,r_q)$ for upper and mirror links, its signed
reflection for lower links, and $[-r_q/2,r_q/2)$ for cut links.
This is a fundamental domain. Reflection maps the lower variables
measure-preservingly to an independent copy of the upper variables,
fixes $(A_0,z_0)$, and changes the sign of $a$.

The upper and mirror photon weights are
\begin{equation}\label{eq:reg-half-photon-weights}
 W_+=\prod_{\sigma\in P_+}e^{-F_\sigma^2/2}
               \prod_{Q\in Q_+}1_{\{(dz)(Q)=0\}},
 \qquad W_0=\prod_{\sigma\in P_0}e^{-F_\sigma^2/2}.
\end{equation}
The full Maxwell weight with the closedness constraint is
$W_0W_+\Theta(W_+)$ times the Gaussian and constraint factors at
the cut plaquettes and cubes. The half weights contain only upper
and mirror variables. Their integrals are bounded uniformly in the
mirror variables by
\[
 \int W_+\,\dd(A_+,z_+)
       \le r_q^{|B_+|}C_q^{|P_+|},\qquad
 \int W_0\,\dd(A_0,z_0)
       \le r_q^{|B_0|}C_q^{|P_0|}.
\]
The proof drops the cube constraints and uses
\eqref{eq:reg-gaussian-sum} on each plaquette.

At a cut cube impose integer closedness by the Fourier multiplier
\[
 1_{\{(dz)(Q)=0\}}
   =\int_0^{2\pi}e^{i\alpha_Q(dz)(Q)}\frac{\dd\alpha_Q}{2\pi}.
\]
For $\alpha\in[0,2\pi)^{Q_\times}$ and
$k\in\Z^{P_\times}$ define
\begin{equation}\label{eq:reg-dual-crossing-variables}
 \begin{aligned}
 \alpha'_\sigma=\sum_{Q\in Q_\times}[Q:\sigma]\alpha_Q,
 \quad \beta_\sigma=\frac{q}{2\pi}(\alpha'_\sigma-2\pi k_\sigma),
 \\
 u_{\alpha,k}
 =\exp\left(i\sum_Q\alpha_Q(d_+z)(Q)
                  -i\sum_\sigma\beta_\sigma(d_+A)_\sigma\right).
 \end{aligned}
\end{equation}
The function $u_{\alpha,k}$ depends only on the upper variables and
has modulus one.

Poisson summation, in the normalization needed here, reads
\begin{equation}\label{eq:reg-poisson-crossing}
 \sum_{j\in\Z}e^{-(y+r_qj)^2/2+ij\alpha}
 =\frac{q}{\sqrt{2\pi}}\sum_{k\in\Z}
          e^{-\beta_k^2/2-i\beta_ky},
 \qquad\beta_k=\frac{q}{2\pi}(\alpha-2\pi k).
\end{equation}
It follows by Fourier transforming the Gaussian under
$t=y+r_qs$. Insert the cube multipliers and apply this formula at
each cut plaquette. Splitting $dA$ and $dz$ by
\eqref{eq:reg-cut-differentials} gives the exact identity
\begin{align}\label{eq:reg-photon-crossing}
 &\sum_{z_\times}
       \prod_{\sigma\in P_\times}e^{-F_\sigma^2/2}
       \prod_{Q\in Q_\times}1_{\{(dz)(Q)=0\}}
 \nonumber\\
 &\quad=\int\prod_{Q\in Q_\times}\frac{\dd\alpha_Q}{2\pi}
       \sum_{k\in\Z^{P_\times}}
       \prod_{\sigma\in P_\times}
               \frac{q}{\sqrt{2\pi}}e^{-\beta_\sigma^2/2}\,
       u_{\alpha,k}\Theta(u_{\alpha,k})
       \exp\left(iq\sum_{b\in B_\times}a_b(d^*k)(b)\right).
\end{align}
Here $k$ is extended by zero outside $P_\times$. To obtain the final
factor, observe that at a cut link $b$,
\[
 \sum_{\sigma\in P_\times}[\sigma:b]\alpha'_\sigma
 =\sum_{Q\in Q_\times}\alpha_Q
                    \sum_\sigma[Q:\sigma][\sigma:b]=0.
\]
The restriction to cut plaquettes may be dropped because any
plaquette containing a cut link is cut; the last sum is zero by
$d^2=0$. Thus the part of $\beta$ involving $\alpha'$ cancels from
the cut-link factor.

All the weights before the two reflected factors on the right of
\eqref{eq:reg-photon-crossing} are positive. Their total integral is
at most $C_q^{|P_\times|}$. Indeed, with $y=0$ in
\eqref{eq:reg-poisson-crossing},
\[
 \frac{q}{\sqrt{2\pi}}\sum_k e^{-\beta_k^2/2}
 \le\sum_j e^{-(r_qj)^2/2}\le C_q
\]
at each cut plaquette. The original Gaussian sum also has an
absolute bound $C_q^{|P_\times|}$, justifying insertion of the
Fourier multipliers and the exchanges of their sums and integrals.

Cut links occur only at coordinate link planes. A link fixed by a
mirror reflection but not contained in its plane has its endpoints
exchanged; hence it is perpendicular to the mirror, whose normal
is a coordinate vector. A diagonal mirror has cut plaquettes and
cubes, but no cut links.

For crossing-mode subsets $I,I'$ let
$\mathfrak q(I\cup I')(b)$ be the sum of their charges at modes lying
over $b$. The fermionic phases in \eqref{eq:reg-two-crossings} are
\[
 u_Iu_{I'}=\exp\left(iq\sum_{b\in B_\times}
                              a_b\mathfrak q(I\cup I')(b)\right).
\]
Integration of the cut links over one period gives
\begin{equation}\label{eq:reg-gauss-indicator}
 \int_{[-r_q/2,r_q/2)^{B_\times}}
 e^{iq\sum_ba_b((d^*k)(b)+\mathfrak q(I\cup I')(b))}\,\dd a
 =r_q^{|B_\times|}
   1_{\{d^*k+\mathfrak q(I\cup I')=0\text{ on }B_\times\}}.
\end{equation}
The quantities in the exponent are integers. This is precisely
Gauss' law at the cut links, with a nonnegative multiplier. When
there are no cut links the condition is empty.

\subsection{The reflected form on an adapted torus}
\label{subsec:reg-reflected-form}

Let $\mathcal B^+$ be the Grassmann-valued integrands whose
generators are in $\mathcal X_+$ and whose bounded measurable
coefficient functions depend only on the upper and mirror
curvatures and link phases. For the later operator estimates we
also use $\mathcal G^+$: its coefficients $\Phi$ satisfy, for some
$C<\infty$ and $\varepsilon>0$,
\begin{equation}\label{eq:reg-upper-growth}
 W_+|\Phi|\le C\prod_{\sigma\in P_+}e^{-\varepsilon F_\sigma^2}.
\end{equation}
This condition is imposed on the upper and mirror variables and
places no growth factor on mirror plaquettes. The bound is relevant
only where the upper cube constraints hold. Such half-configurations
can be extended to full closed configurations by reflecting their
upper $z$, keeping $z_0$, and taking $z_\times=0$; the cut-cube
identity \eqref{eq:reg-cut-differentials} proves closedness. Thus
these coefficient functions have an unambiguous reading wherever
$W_+\ne0$. In particular $\mathcal B^+\subset\mathcal G^+$.

\begin{proposition}\label{thm:reg-adapted-positivity}
For a torus adapted to $\ell$ and $G,G'\in\mathcal G^+$, the integral
$Z_{\mathbb T}[\Theta(G')G]$ converges absolutely, coefficient by
coefficient, and admits a representation
\begin{equation}\label{eq:reg-positive-representation}
 Z_{\mathbb T}[\Theta(G')G]
   =\int\overline{\widehat{G'}(s)}\widehat G(s)\,\dd\mathfrak m(s),
\end{equation}
where $\mathfrak m$ is a finite positive measure and
$G\mapsto\widehat G$ is linear with bounded measurable values.
This holds for even, odd and mixed-parity integrands. In particular
$Z_{\mathbb T}[\Theta(G)G]\ge0$.

For every bounded scalar observable $\Phi$ depending only on curvatures
of plaquettes wholly contained in $U_+$, the curvature law satisfies
\begin{equation}\label{eq:reg-law-positivity}
 \int_\Omega\overline{\Phi(\theta^*F)}\Phi(F)\,\dd\mu_{\mathbb T}(F)
 \ge0.
\end{equation}
It is invariant under this reflection.
\end{proposition}

\begin{proof}
The measure space has coordinates
\[
 s=(A_0,z_0,\alpha,k,I,I'),
\]
with mirror variables $(A_0,z_0)$, cut-cube variables
$\alpha\in[0,2\pi)^{Q_\times}$, integers $k\in\Z^{P_\times}$, and
the two finite mode subsets. Define
\begin{align}\label{eq:reg-positive-measure}
 \dd\mathfrak m(s)={}&W_0\,\dd(A_0,z_0)
     \prod_{Q\in Q_\times}\frac{\dd\alpha_Q}{2\pi}
     \prod_{\sigma\in P_\times}
                         \frac{q}{\sqrt{2\pi}}e^{-\beta_\sigma^2/2}
 \nonumber\\
 &\quad\cdot\omega_I\omega_{I'}r_q^{|B_\times|}
        1_{\{d^*k+\mathfrak q(I\cup I')=0\text{ on }B_\times\}},
\end{align}
where $z_0,k,I,I'$ carry counting measure. Put
\begin{equation}\label{eq:reg-half-transform}
 \widehat G(s)=\int W_+u_{\alpha,k}
         \int_{\mathcal X_+}G e^{-S_+}\zeta_I\check\zeta_{I'}
                                      \,\dd(A_+,z_+).
\end{equation}
The action $S_+$ has bounded coefficients depending only on upper
and mirror links. Thus the integrand in this formula is a
half-integrand of the stated kind.

We first prove the identity for bounded coefficients and definite
parities. Inserting \eqref{eq:reg-two-crossings} into the reflected
functional combines each term with the half factors
\[
 \xi=G e^{-S_+}\zeta_I\check\zeta_{I'},\qquad
 \xi'=G'e^{-S_+}\zeta_I\check\zeta_{I'}.
\]
The sign from this rearrangement is
$(-1)^{(|G|+|G'|)(|I|+|I'|)}$. It is one when the parities agree.
When they disagree, one of $\xi,\xi'$ is odd, so its half Berezin
integral vanishes; the corresponding full integral also vanishes
by \eqref{eq:reg-half-berezin}. The sign consequently disappears
from the integrated formula in all cases.

Now insert \eqref{eq:reg-photon-crossing} and integrate the cut
links. Formula \eqref{eq:reg-gauss-indicator} supplies the last
factor in \eqref{eq:reg-positive-measure}. At fixed $s$, the
remaining factors are
\[
 \Theta\bigl(W_+u_{\alpha,k}\xi'\bigr)
                    W_+u_{\alpha,k}\xi.
\]
The lower photon variables map measure-preservingly onto an
independent set of upper variables. Combining this change of
variables with \eqref{eq:reg-half-berezin} gives
$\overline{\widehat{G'}(s)}\widehat G(s)$, proving
\eqref{eq:reg-positive-representation} once the exchanges of
integration have been justified.

For that justification let $C$ bound all the finitely many half
Berezin coefficients of $G e^{-S_+}\zeta_I\check\zeta_{I'}$ and of
the corresponding expression with $G'$. The scalar phase factors
have modulus one. Before the cut-link integration, an absolute
majorant therefore has total integral at most
\begin{equation}\label{eq:reg-representation-majorant}
 C^2(r_qC_q)^{3V}
       \prod_{i\text{ at }\Pi}(1+\omega_i)
       \prod_{i'\text{ at }\Pi'}(1+\omega_{i'}).
\end{equation}
The upper, lower and mirror variables are bounded by their
Gaussian sums; the cut-link intervals contribute their lengths;
the dual photon weights are bounded by $C_q^{|P_\times|}$; and the
finite mode sums give the two products. This proves absolute
integrability and all required uses of Fubini's theorem. The
same bounds without the half factors show that $\mathfrak m$ is
finite. It is positive by construction. Formula
\eqref{eq:reg-half-transform} is measurable and bounded uniformly
in $s$ by $Cr_q^{|B_+|}C_q^{|P_+|}$. Mixed parities follow by
sesquilinearity.

For $G,G'\in\mathcal G^+$, use
\eqref{eq:reg-upper-growth} with exponents $\varepsilon$ and
$\varepsilon'$. The algebraic identities above are unchanged.
The absolute majorant now replaces the upper and lower Gaussian
sums by $C_q(\varepsilon)^{|P_+|}$ and
$C_q(\varepsilon')^{|P_-|}$, respectively, while the mirror and
cut sums retain exponent $1/2$. These quantities are finite by
\eqref{eq:reg-gaussian-sum}. This proves coefficientwise
integrability, the same representation, and the boundedness of
$\widehat G$. Taking $G'=G$ proves positivity.

For a scalar $\Phi$ with this plaquette dependence, divide this positivity
by $Z_{\mathbb T}[1]>0$ and use the definition of the curvature
law. Reflection invariance follows from
\eqref{eq:reg-reflected-functional}, since the adapted torus is
preserved by reflection. This proves
\eqref{eq:reg-law-positivity} and completes the theorem.
\end{proof}

The same form remains positive after multiplication of the weight
by $\Phi_+\Theta(\Phi_+)$, with $\Phi_+$ bounded and scalar on the
upper half: apply the theorem to $G\Phi_+$. This observation will
be used for the reflected arrays of observables. It uses positivity
of the joint reflected form before integrating out the fermion.
The determinant's pointwise nonnegativity by itself does not imply
reflection positivity.


